\chapter{Felhasználói dokumentáció}\label{ch:problem-definition}

\section{Felhasznált módszerek röviden}

\todo{Megírnit hogyan epul fel az app roviden}

\section{Rendszerkövetelmények}

A szerver futattasához szüksegés egy olyan környezet, amin futtatható a Docker\footnote{\url{https://www.docker.com/}. Utolsó elérés: 2026.04.10.}. Emiatt az alabbi követelmények szükségesek:
\begin{itemize}
	\item 64 bit-es Windows 11 vagy Linux operációs rendszer ami támogatja a Docker futtatását
	\item Legalább 4 GB RAM, ajánlott 8 GB 
	\item Legalább 10 GB szabad tárhely
	\item Hálózati kapcsolat
\end{itemize}

Kliens oldalon egy modern webböngészőre van szükségünk, ami támogatja a HTML5-öt és a JavaScript-et:
\begin{itemize}
	\item Mozilla Firefox legfrissebb verziója
	\item Google Chrome legfrissebb verziója
	\item Microsoft Edge legfrissebb verziója
	\item Apple Safari legfrissebb verziója
\end{itemize}

\section{Telepítés}
Az alkalmazást a korábban említett Docker\footnote{\url{https://docs.docker.com/desktop/}. Utolsó elérés: 2026.04.10.}, illetve Docker Compose\footnote{\url{https://docs.docker.com/compose/install/}. Utolsó elérés: 2026.04.10.} segítségével lehet telepíteni. Ennek köszönhetően egyetlen parancs kiadásával a teljes rendszer, különösebb beállítások nélkül elindítható.

\subsection{Telepítési útmutató}
Miután a Docker és a Docker Compose telepítésre került, fájkezelőben a projekt mappájába kell navigálni. Azt itt található \texttt{docker-compose.yml} fájl segítségével indítható el a rendszert. Ez a fájl négy konténerből áll: egy PostgreSQL\footnote{\url{https://www.postgresql.org/}. Utolsó elérés: 2026.04.10.} adatbázisból, a backend REST API-ból, a frontend szerverből és egy Nginx\footnote{\url{https://nginx.org/}. Utolsó elérés: 2026.04.10.} fordított proxyból.

A rendszer konfigurációja egy environment fájlon keresztül történik, amelyben a működéshez szükséges környezeti változók adhatók meg. Az elérhető konfigurációs opciókat az \ref{tab:config}~táblázat tartalmazza.

Elindításához Linux és macOS rendszereken a \texttt{start.sh}, Windows rendszereken pedig a \texttt{start.ps1} scriptet kell futtatni, amely automatikusan elindítja a szükséges konténereket a konfigurációs fájl alapján. 

Az alap beállításokkal a rendszer a \texttt{http://localhost:9393} címen lesz elérhető, amit egy webböngészőben lehet megnyitni.

\section{Funkciók}
Az alkalmazásban két fajta felhasználói szerepkör van: adminisztrátor és normál felhasználó. Az oldalon a regisztráció nélkül böngészhetően a receptek. Bejelentkezés után a normál felhasználók hozzáférhetnek a közösségi funkciókhoz, kezelhetik saját receptjeiket, míg az adminisztrátorok további jogosultságokkal rendelkeznek, például a felhasználók kezeléséhez.

\subsection{Receptkezelés}
A felület könnyen lehetővé teszi a felhasználók számára, hogy létrehozzanak, szerkeszthessenek és törölhessenek recepteket. Minden recepthez megadható egy cím, leírás, hozzávalók listája és elkészítési útmutató. Egy képet is feltölthetnek a recepthez. 

\subsection{Közösségi funkciók}

A regisztrált felhasználók értékelhetik és kommentelhetik a recepteket. Az értékelése egytől ötig terjedő skálán történik, amelyek csillag formájában jellennek meg. A receptekhez kommenteket is lehet írni, amelyeket a receptek alatt lehet megtekinteni. 
\subsection{Keresés}

Az alkalmazás lehetőséget biztosít a receptek közötti keresésre. A keresőmezőbe kulcsszavakat lehet beírni, amelyek alapján a rendszer megjeleníti a releváns recepteket. Emellett a receptek szűkíthetőek allergének alapján, így a felhasználók könnyedén megtalálhatjak a számukra megfelelő recepteket.

\subsection{Profilkezelés}

A felhasználók saját profiljukat is kezelhetik, ahol megtekinthetik és szerkeszthetik profilképüket, adataikat, valamint megváltoztathatják jelszavukat. A profilban láthatóak az  általuk létrehozott receptek is, amelyeket könnyen elérhetnek és szerkeszthetnek.

\subsection{Adminisztrátor funkciók}

Az adminisztrátorok hozzáférnek az oldalon több funkcióhoz, mint a normál felhasználók. Az adminisztrátorok kezelhetik a felhasználókat, például törölhetnek vagy módosíthatnak fiókokat. Emellett az adminisztrátorok moderálhatják a recepteket és a kommenteket, biztosítva ezzel a közösségi tartalom minőségét.

